% ----------------------------------------------------------------
% LaTeX definitions for a Beamer presentation     ****************
% Template prepared by T.A. McWalter (April 2010)
% ----------------------------------------------------------------
\documentclass{beamer}
\let\Tiny=\tiny

% BEAMER OPTIONS -------------------------------------------------

\mode<presentation>
{
  \useinnertheme[shadow]{rounded}
  \usetheme{Copenhagen}
  \setbeamercovered{transparent}
  \useoutertheme{infolines}
  \useoutertheme{shadow}
%  \setbeamertemplate{navigation symbols}{}     % Switches off navigation icons
}

% PACKAGES -------------------------------------------------------

\usepackage{amsmath}
\usepackage{amsbsy}
\usepackage{amsfonts}
\usepackage{amssymb}
\usepackage{color}
\usepackage{ifpdf}
\usepackage{mathrsfs}
\usepackage{multirow}
\usepackage{natbib}
\usepackage{xspace}

% GRAPHICS HANDLING ----------------------------------------------
\ifpdf
    \usepackage{graphicx}
    \DeclareGraphicsExtensions{.pdf}
    \newcommand{\bb}{viewport}
\else
    \usepackage{graphicx}
    \DeclareGraphicsExtensions{.eps}
    \newcommand{\bb}{bb}
\fi

%\usepackage{pgf}
%\usepackage{pgfpages}
%\pgfpagesuselayout{resize to}[a4paper,border shrink=5mm,landscape]     % Use this option for printing on A4 paper
%\pgfpagesuselayout{4 on 1}[a4paper,border shrink=5mm,landscape]		% Use this option for printing 4 slides per page on A4 paper

% THEOREM STYLES -------------------------------------------------

\newtheorem{thm}{Theorem.}
\newtheorem{cor}[thm]{Corollary.}
\newtheorem{con}[thm]{Conjecture.}
\newtheorem{ass}[thm]{Assumption.}
\newtheorem{lem}[thm]{Lemma.}
\newtheorem{prop}[thm]{Proposition.}
\newtheorem{defn}[thm]{Definition.}
\renewcommand{\thethm}{}  % not numbered

% PREAMBLE -------------------------------------------------------

\title[Convertible Bonds with Credit Risk]{Valuation of Convertible Bonds with Credit Risk}

%\subtitle{Subtitle (if needed)}

\author[D.B. Naylor]{David Naylor}

\institute[AMF Wits]{Programme in Advanced Mathematics of Finance,\\
    School of Computational and Applied Mathematics,\\
    University of the Witwatersrand\\
    \vspace{.5cm}\includegraphics[height=2cm]{../common/WitsColourLogo}}

\date{19 November 2012}

\def\n{\nonumber\\}
\def\h{\\\hline}
\newcommand{\theader}[1]{\multicolumn{1}{|c|}{\textbf{#1}}}

% Bond descriptions
\def\t{\ensuremath{t}\xspace}
\def\ti{\ensuremath{t_i}\xspace}
\def\T{\ensuremath{T}\xspace}
\def\N{\ensuremath{N}\xspace}
\def\R{\ensuremath{R}\xspace}
\def\C{\ensuremath{C_i}\xspace}
\def\Kp{\ensuremath{K^{p}_t}\xspace}
\def\Kc{\ensuremath{K^{c}_t}\xspace}
\def\Kv{\ensuremath{\kappa_t}\xspace}
\def\Tp{\ensuremath{\Omega^{p}}\xspace}
\def\Tc{\ensuremath{\Omega^{c}}\xspace}
\def\Tv{\ensuremath{\Omega^{v}}\xspace}
\def\ST{\ensuremath{S_T}\xspace}

% Statistical
\def\E[#1]{\Tilde{\mathbb{E}}\!\left[#1\right]}
\def\Et[#1]{\Tilde{\mathbb{E}}_t\!\left[#1\right]}
\def\Var[#1]{\Tilde{\mathbb{V}}\mathtt{ar}\!\left[#1\right]}
\def\Prob(#1){\Tilde{\mathbb{P}}\!\left(#1\right)}

% Derivatives
\def\S{\ensuremath{S_t}\xspace}
\def\Sdt{\ensuremath{S_{t + \dt}}\xspace}
\def\SDt{\ensuremath{S_{t + \Dt}}\xspace}
\def\W{\ensuremath{\tilde{W}_t}\xspace}
\def\q{\ensuremath{q_t}\xspace}
\def\qdt{\ensuremath{q_{t + \dt}}\xspace}
\def\dP{\ensuremath{d\Pi_t}\xspace}
\def\dS{\ensuremath{dS_t}\xspace}
\def\dt{\ensuremath{dt}\xspace}
\def\dW{\ensuremath{d\W}\xspace}
\def\dq{\ensuremath{d\q}\xspace}
\def\Vt{\frac{\partial V}{\partial t}}
\def\Vs{\frac{\partial V}{\partial S}}
\def\Vss{\frac{\partial^2 V}{\partial S^2}}

% Differences
\def\DS{\ensuremath{\delta S_t}\xspace}
\def\Dt{\ensuremath{\delta t}\xspace}
\def\DW{\ensuremath{\delta \W}\xspace}
\def\Dq{\ensuremath{\delta \q}\xspace}

% Finite difference
\def\Vjn{\ensuremath{V^{n}_j}\xspace}
\def\Vjnm{\ensuremath{V^{n - 1}_j}\xspace}
\def\Vjnp{\ensuremath{V^{n + 1}_j}\xspace}
\def\Vjmn{\ensuremath{V^{n}_{j - 1}}\xspace}
\def\Vjpn{\ensuremath{V^{n}_{j + 1}}\xspace}
\def\Vjmnm{\ensuremath{V^{n - 1}_{j - 1}}\xspace}
\def\Vjpnm{\ensuremath{V^{n - 1}_{j + 1}}\xspace}
\def\Vjmnp{\ensuremath{V^{n + 1}_{j - 1}}\xspace}
\def\Vjpnp{\ensuremath{V^{n + 1}_{j + 1}}\xspace}

\def\Von{\ensuremath{V^{n}_{0}\xspace}}
\def\Vomn{\ensuremath{V^{n}_{-1}\xspace}}
\def\Vopn{\ensuremath{V^{n}_{1}\xspace}}
\def\VJn{\ensuremath{V^{n}_{J}\xspace}}
\def\VJmn{\ensuremath{V^{n}_{J - 1}\xspace}}
\def\VJpn{\ensuremath{V^{n}_{J + 1}\xspace}}
\def\Vonm{\ensuremath{V^{n - 1}_{0}\xspace}}
\def\Vomnm{\ensuremath{V^{n - 1}_{-1}\xspace}}
\def\Vopnm{\ensuremath{V^{n - 1}_{1}\xspace}}
\def\VJnm{\ensuremath{V^{n - 1}_{J}\xspace}}
\def\VJmnm{\ensuremath{V^{n - 1}_{J - 1}\xspace}}
\def\VJpnm{\ensuremath{V^{n - 1}_{J + 1}\xspace}}

\def\Sz{\ensuremath{S_0}\xspace}
\def\Sj{\ensuremath{S_j}\xspace}
\def\SJ{\ensuremath{S_J}\xspace}
\def\Xj{\ensuremath{X_j}\xspace}

% Binomial model
\def\u{\ensuremath{u}\xspace}
\def\d{\ensuremath{d}\xspace}
\def\o{\ensuremath{o}\xspace}

\def\pu{\ensuremath{p_\u}\xspace}
\def\pd{\ensuremath{p_\d}\xspace}
\def\po{\ensuremath{p_\o}\xspace}

% Binomial model valuation
\def\V{\ensuremath{V_t}\xspace}
\def\Vu{\ensuremath{V^{u}_t}\xspace}
\def\Vd{\ensuremath{V^{d}_t}\xspace}
\def\X{\ensuremath{X_t}\xspace}
\def\P{\ensuremath{\Pi_t}\xspace}
\def\PDt{\ensuremath{\Pi_{t + \Dt}}\xspace}
\def\D{\ensuremath{\Delta_t}\xspace}

% Sundry
\def\th{\textsuperscript{th}\xspace}
\def\half{\tfrac{1}{2}}


% =============================================================================

\begin{document}

\begin{frame}
  \titlepage
\end{frame}

% =============================================================================

\section{Introduction}

\begin{frame}{Introduction}
Project Objective:
\begin{itemize}
 \item Understand and implement pricing of a convertible bond with credit risk:
       \begin{itemize}
        \item Finite Difference Model
        \item Binomial Model
       \end{itemize}
\end{itemize}
\end{frame}

\begin{frame}{Literature}
Finite Difference Model:
\begin{itemize}
 \item \Citet{HP96} and \citet{TF98} modelled the derivative using geometric Brownian motion
 \item \Citet{AFV03} modelled non-total default jump of the stock
\end{itemize}

Binomial Model:
\begin{itemize}
 \item \Citet{DL02} first proposed a trinomial tree
 \item \Citet{BEA94}, and \citet{H11} have implemented a binomial tree
 \item \Citet{MK12} implemented a binomial model that converges to the stochastic model
\end{itemize}
\end{frame}

% =============================================================================

\begin{frame}{Outline of the Talk}
  \tableofcontents[hidesubsections]
\end{frame}

% =============================================================================

\AtBeginSection[] {
    \begin{frame}<beamer>{Outline}
        \tableofcontents[currentsection,hideothersubsections]
    \end{frame}
}

% =============================================================================

\section{Convertible Bond} % 7 mins

% -----------------------------------------------------------------------------

\subsection{Convertible Bond definition}

\begin{frame}{Convertible Bond (1/3)}
\begin{defn}[Convertible Bond]
A convertible bond has the following components:
\begin{itemize}
 \item \emph{Annuity:} a series of coupons with a final redemption payment
 \item \emph{Put:} a put option on the bond for the \underline{bond holder}
 \item \emph{Call:} a call option on the bond for the \underline{bond issuer}
 \item \emph{Conversion:} an asset swap for stock by the \underline{bond holder}
\end{itemize}
\end{defn}
\end{frame}

\begin{frame}{Convertible Bond (2/3)}
\begin{defn}[Convertible Bond]
A convertible bond has the following components\footnote{\V denoted the intrinsic value of the derivative}:
\begin{itemize}
 \item \emph{Annuity:} Coupons: \C at \ti; Redemption: \R at \T
 \item \emph{Put:} Strike: \Kp for $\t \in \Tp$; Payoff: $\Kp 1_{(\V \le \Kp)\wedge(\t \in \Tp)}(\V, \t)$
 \item \emph{Call:} Strike: \Kc for $\t \in \Tc$; Payoff: $\Kc 1_{(\V \ge \Kc)\wedge(\t \in \Tc)}(\V, \t)$
 \item \emph{Conversion:} Stock: $\Kv\S$ for $\t \in \Tv$; Payoff: $\Kv\S 1_{(\V \le \Kv\S)\wedge(\t \in \Tv)}(\V, \t)$
\end{itemize}
\end{defn}
\end{frame}

\begin{frame}{Convertible Bond (3/3)}
\begin{defn}[Convertible Bond]
A convertible bond has the following components:
\begin{itemize}
 \item \emph{Annuity:} $\ti \not= \T$
 \item \emph{Put:} $\T \not\in \Tp$
 \item \emph{Call:} $\T \not\in \Tv$; $\Kc > \Kp$
 \item \emph{Conversion:} $\kappa > 0$; $\T \in \Tv$; supersedes the put and call option
\end{itemize}
\end{defn}
\end{frame}

\begin{frame}{Payoff and Actions}
\begin{table}[ht]
 \footnotesize
 \centering
 \begin{tabular}{|l|l|l|}\hline
  \theader{Action}    & \theader{Payoff} & \theader{Condition} \h
  Put                 & \Kp              & $(\V \le \Kp)\wedge(\t \in \Tp)\wedge[(\Kv\S < \V)\lor(\t \not\in \Tv)]$ \h
  Call                & \Kc              & $(\V \ge \Kc)\wedge(\t \in \Tc)\wedge[(\Kv\S < \Kc)\lor(\t \not\in\Tv)]$ \h
  Conversion          & $\Kv\S$          & $(\Kv\S \ge \V)\wedge(\t \in \Tv)$ \h
  Forced conversion   & $\Kv\S$          & $(\V > \Kv\S \ge \Kc)\wedge(\t \in \Tv)\wedge(\t \in \Tc)$ \h
  Redemption          & \R               & $(\t = \T)\wedge[(\Kv\S \le \R)\lor(\t \not\in \Tv)]$ \h
  Hold                &                  & \emph{otherwise} \h
 \end{tabular}
 \caption[Payoff for the convertible bond]{Payoff for the convertible bond.  \V is the intrinsic value of the derivative}
 \label{tab:choices}
\end{table}
\end{frame}

\begin{frame}{Payoff and Actions}
\begin{figure}[ht]
 \centering
 \includegraphics[width=10cm]{../common/fig_choices}
\end{figure}
\end{frame}


% -----------------------------------------------------------------------------

\subsection{Credit Risk}

\begin{frame}{Credit Risk}
\begin{defn}[Convertible Bond on Default]
In default the convertible bond has the following components:
\begin{itemize}
 \item \emph{Annuity:} residual value\footnote{This is assuming the bond recovery is based on the redemption value.  Another possibility is to base the recovery on the bond value at time \t.} of $\gamma\R$
 \item \emph{Conversion:} stock price of $(1 - \eta)\Kv\S$
\end{itemize}

The value (payoff) the convertible bond is:
\begin{equation}
 \X = \max(\gamma\R, (1 - \eta)\Kv\S)
\end{equation}

Default is considered terminal.
\end{defn}
\end{frame}

% =============================================================================

\section{Finite Difference Model} % 4 mins

% -----------------------------------------------------------------------------

\subsection{Stochastic Process}

\begin{frame}{Stock Process with Credit Risk (1/3)}
\begin{defn}[Geometric Brownian with Credit Risk]
A stock process with geometric Brownian motion, with a jump component to model the default event, has the following stochastic differential equation:
\begin{equation}
 \dS = \mu\S\dt + \sigma\S\dW - \eta\S\dq
\end{equation}
with \q being a Poisson process with intensity $\lambda$ where the first event $\q = 1$ is the default event.
\end{defn}
\end{frame}

\begin{frame}{Stock Process with Credit Risk (2/3)}
\begin{cor}
The stochastic process \q, when representing terminal default, follows an exponential distribution with:
\begin{align}
 \q                       &\sim \exp(\lambda) \\
 \Prob(\qdt = 1 | \q = 0) &= \lambda\dt \\
 \E[\dq]                  &= \lambda\dt \\
 \Var[\dq]                &= \lambda\dt
\end{align}
\end{cor}
\end{frame}

\begin{frame}{Stock Process with Credit Risk (3/3)}
\begin{cor}
Under risk neutral probability space the drift rate of \S is $\mu = (r + \lambda\eta)$ and the stochastic differential equation is:
\begin{equation}
 \dS = (r + \lambda\eta)\S\dt + \sigma\S\dW - \eta\S\dq
\end{equation}
\end{cor}
\end{frame}

% -----------------------------------------------------------------------------

\subsection{Partial Differential Equation}

\begin{frame}{Partial Differential Equation}
\begin{thm}[Black-Scholes with Default]
The price of an option, on an underlying \S, has the following partial differential equation:
\begin{equation}
 (r + \lambda)\V = \Vt + (r + \lambda\eta)\S\Vs + \half\sigma^2\S^2\Vss + \lambda\X \label{eq:BScr}
\end{equation}
\end{thm}
\end{frame}

\begin{frame}{Derivation}
The following steps are used in the derivation for Equation \ref{eq:BScr}:
\begin{enumerate}
 \item Create a portfolio $\P = \V - \D\S$
 \item Invest the residual of \P at the risk free rate: $\dP = r\P\dt$
 \item Equate the expected value of the derivative of \P with the risk free rate:
       \begin{align}
        r\P\dt = \E[\dP] &= \Prob(\qdt = 0) \E[d(\V - \D\S)|\qdt = 0] + \n
                         &\qquad \Prob(\qdt = 1) \E[d(\V - \D\S)|\qdt = 1]
       \end{align}
 \item Note that the derivative of \P when going into default is:
       \begin{equation}
        \E[d(\V - \D\S)|\qdt = 1] = \X - (1 - \eta)\D\S -(\V - \D\S) \label{eq:dPx}
       \end{equation}
 \item Choose $\D = \Vs$ to eliminate the \dW terms
\end{enumerate}
\end{frame}


% =============================================================================

\section{Binomial Model} % 9 mins

% -----------------------------------------------------------------------------

\subsection{Binomial Tree}

\begin{frame}{Binomial Tree (1/2)}
\begin{defn}[Binomial Process with Default]
A binomial process with up movement \u and down movement \d and a terminal default event $(1 - \eta)$ has the following mass distribution function:
\begin{equation}
 f(\SDt) = \left\{\begin{array}{l@{\qquad\mbox{if }}l} \pu & \SDt = u\S \\ \pd & \SDt = d\S \\ \po & \SDt = (1 - \eta)\S \end{array} \right.
\end{equation}
\end{defn}
\end{frame}

\begin{frame}{Binomial Tree (2/2)}
\begin{cor}
The expected value and variance of the Binomial Process with Default is:
\begin{align}
 \E[\frac{\SDt}{\S}]   &= \u\pu + \d\pd + (1 - \eta)\po \\
 \Var[\frac{\SDt}{\S}] &= \u^2\pu + \d^2\pd + (1 - \eta)^2\po - \E[\frac{\SDt}{\S}]^2
\end{align}
\end{cor}
\end{frame}

% -----------------------------------------------------------------------------

\subsection{Binomial Process}

\begin{frame}{Binomial Process (1/2)}
\begin{thm}[Binomial Process with Default]
A Binomial Process with Default has the following parameters:
\begin{align*}
  \u  &= e^{\sigma\sqrt{\Dt}} \\
  \d  &= e^{-\sigma\sqrt{\Dt}} \\
  \pu &= \frac{e^{r\Dt} - \d e^{-\lambda\Dt} - (1 - \eta)(1 - e^{-\lambda\Dt})}{\u - \d} \\
  \pd &= e^{-\lambda\Dt} - \pu \\
  \po &= 1 - e^{-\lambda\Dt} \\
  \V  &= e^{-r\Dt}(\Vu\pu + \Vd\pd + \X\po) + c_i \\
\end{align*}
\end{thm}
\end{frame}

\begin{frame}{Binomial Process (2/2)}
\begin{thm}[Binomial Process with Default]
A Binomial Process with Default has the following restrictions on the parameters:
\begin{align*}
  0   &< \sigma \\
  0   &\le \lambda \\
  0   &< \Dt \\
  \Dt &\le \frac{\sigma^2}{r^2} \\
  \Dt &\le \frac{1}{\lambda}\ln\left(\frac{u - (1 - \eta)}{e^{r\Dt} - (1 - \eta)}\right) \\
\end{align*}
\end{thm}
\end{frame}

\begin{frame}{Derivation (1/3)}
The following steps are used in the derivation for the parameters of the Binomial Process with Default.
\begin{enumerate}
 \item Note that the arrival time of a default event follows an exponential distribution with hazard rate $\lambda$:
       \begin{align}
        \po &= 1 - e^{-\lambda\Dt} \\
        \pd &= e^{-\lambda\Dt} - \pu
       \end{align}
 \item Equate $\E[\frac{\SDt}{\S}]$ with $\E[\frac{\Sdt}{\S}] = e^{r\Dt}$, where $\dt \approxeq \Dt$
 \item Equate $\Var[\frac{\SDt}{\S}]$ with $\Var[\frac{\Sdt}{\S}] = (\sigma^2 + \lambda\eta^2)\Dt$, where $\dt \approxeq \Dt$
       \begin{enumerate}
        \item Assume $\Dt^2$ = 0
        \item $\u\d = 1$
        \item $\u = e^{\sqrt{A\Dt}}$
        \item Use Taylor series expansion for all exponential terms
       \end{enumerate}
\end{enumerate}
\end{frame}

\begin{frame}{Derivation (2/3)}
The following steps are used in the derivation for the valuation of the Binomial Process with Default.
\begin{enumerate}
 \item Construct a portfolio $\P = \V - \D\S$ with:
       \begin{equation}
        \PDt = \left\{
                \begin{array}{l@{\qquad\mbox{with probability }}l}
                 \Vu - \D\S\u      & \pu \\
                 \Vd - \D\S\d      & \pd \\
                 \X - (1 - \eta)\S & \po
                \end{array}\right.
       \end{equation}
 \item Invest the residual of \P at the risk free rate: $\PDt = \P e^{r\Dt}$
 \item Equate the expected value of \PDt with the risk free rate:
       \begin{align}
        \P e^{r\Dt} = \E[\PDt]
       \end{align}
 \item Choose $\D = \frac{\Vu - \Vd}{\S(u - d)}$ to hedge against up and down movements of the stock
\end{enumerate}
\end{frame}

\begin{frame}{Derivation (3/3)}
The following steps are used in the derivation for the limits of the Binomial Process with Default.
\begin{enumerate}
 \item Note that $\{\pu, \pd, \po\}$ is required to be a valid probability set:
       \begin{align}
        \min(\pu, \pd, \po) \ge 0 \\
        \pu + \pd + \po = 1
       \end{align}
 \item Using the relationship between \pu, \pd and \po the following inequality needs to hold:
       \begin{equation}
        0 \le \pu \le e^{-\lambda\Dt}
       \end{equation}
 \item Impose the requirements that $\lambda\Dt \ge 0$
\end{enumerate}
\end{frame}

% =============================================================================

\section{Numerical Example}

\begin{frame}{Parameters for the Convertible Bond}
Consider the following parameters of a Convertible Bond:
\begin{table}[ht]
 \centering
 \begin{tabular}{|c|l|c|} \hline
  \theader{Component}         & \theader{Parameter} & \theader{Value} \h
  \multirow{5}{*}{Annuity}    & Notional            & $100$ \\
                              & Coupon              & 8\% \\
                              & Coupon frequency    & Semi-annually \\
                              & Maturity            & $\T := 5$ \\
                              & Recovery            & $\gamma := 0\%$ \h
  \multirow{2}{*}{Put}        & Strike              & $\Kp := 105$ \\
                              & Period(s)           & $\Tp := \{3\}$ \h
  \multirow{2}{*}{Call}       & Strike              & $\Kc := 110 + \C\tfrac{t \pmod {0.5}}{0.5}$ \\
                              & Period(s)           & $\Tc := [2, 5)$ \h
  \multirow{2}{*}{Conversion} & Quantity of stocks  & $\Kv := 1$ \\
                              & Period(s)           & $\Tv := [0, 5]$ \h
 \end{tabular}
 \caption{Convertible Bond Parameters}
\end{table}
\end{frame}

\begin{frame}{Parameters for the Stock}
\begin{table}[ht]
 \centering
 \begin{tabular}{|l|l|} \hline
  \theader{Parameter} & \theader{Value} \h
  Risk free rate      & $r := 5\%$ \h
  Volatility          & $\sigma := 20\%$ \h
  Hazard rate         & $\lambda := 2\%$ \h
  Default (total)     & $\eta := 100\%$ \h
  Default (typical\footnote{\citet{BP95} found that stock prices typically drop 30\% on announcement of default}) & $\eta := 30\%$ \h
  Default (partial)   & $\eta := 0\%$ \h
 \end{tabular}
 \caption{Stock Parameters}
 \label{tab:stock_parameters}
\end{table}
\end{frame}

\begin{frame}{Payoff and Action}
\begin{figure}[ht]
 \centering
 \includegraphics[width=8cm]{../common/fig_choices}
 \caption[Payoff]{Payoff of the convertible bond with total default and colours indicating the action taken for that payoff.  $\Dt = 2^{-3}$, $\DS = 2^{1}$ and $\S \in [0, 250]$.}
\end{figure}
\end{frame}

\begin{frame}{Simplified Call}
\begin{figure}[ht]
 \centering
 \includegraphics[width=12cm]{../common/fig_simple}
 \caption[Simple Call]{Comparison of initial value of the standard convertible bond compared to one with a singular time for the call provision.  Also the payoff of the ``simple call'' and colours indicating the action taken for the payoff.}
\end{figure}
\end{frame}

\begin{frame}{Varying Call Time}
\begin{figure}[ht]
 \centering
 \includegraphics[width=8cm]{../common/fig_varTc}
 \caption[Varying Call Time]{Comparison of initial value of convertible bonds with different times for the call provision.}
\end{figure}
\end{frame}

\begin{frame}{Varying Put Time}
\begin{figure}[ht]
 \centering
 \includegraphics[width=8cm]{../common/fig_varTp}
 \caption[Varying Put Time]{Comparison of initial value of convertible bonds with different times for the put provision.}
\end{figure}
\end{frame}

\begin{frame}{Varying Redemption Value}
\begin{figure}[ht]
 \centering
 \includegraphics[width=12cm]{../common/fig_varR}
 \caption[Varying Redemption Value]{Comparison of initial value of standard convertible bonds to one with no redemption value.  Also the payoff of the ``no redemption'' and colours indicating the action taken for the payoff.}
\end{figure}
\end{frame}
\section{Bibliography}

\begin{frame}{Bibliography}
\scriptsize
\bibliographystyle{../references/amsplnat}
\bibliography{../references/references.bib}
\end{frame}
\end{document}
